\begin{proposition}
\label{proposition-perfect-is-compact}
Let $R$ be a ring. For an object $K$ of $D(R)$ the following are equivalent
\begin{enumerate}
\item $K$ is perfect, and
\item $K$ is a compact object of $D(R)$.
\end{enumerate}
\end{proposition}

\begin{proof}
Assume $K$ is perfect, i.e., $K$ is quasi-isomorphic to a bounded complex
$P^\bullet$ of finite projective modules, see
Definition \ref{definition-perfect}. If $E_i$ is represented by the complex
$E_i^\bullet$, then $\bigoplus E_i$ is represented by the complex
whose degree $n$ term is $\bigoplus E_i^n$. On the other hand,
as $P^n$ is projective for all $n$ we have
$\Hom_{D(R)}(P^\bullet, K^\bullet) = \Hom_{K(R)}(P^\bullet, K^\bullet)$
for every complex of $R$-modules $K^\bullet$, see
Derived Categories,
Lemma \ref{derived-lemma-morphisms-from-projective-complex}.
Thus $\Hom_{D(R)}(P^\bullet, E^\bullet)$ is the cohomology of the complex
$$
\prod \Hom_R(P^n, E^{n - 1}) \to
\prod \Hom_R(P^n, E^n) \to
\prod \Hom_R(P^n, E^{n + 1}).
$$
Since $P^\bullet$ is bounded we see that we may replace the $\prod$
signs by $\bigoplus$ signs in the complex above. Since each $P^n$ is a finite
$R$-module we see that
$\Hom_R(P^n, \bigoplus_i E_i^m) = \bigoplus_i \Hom_R(P^n, E_i^m)$
for all $n, m$.
Combining these remarks we see that the map of
Derived Categories, Definition \ref{derived-definition-compact-object}
is a bijection.

\medskip\noindent
Conversely, assume $K$ is compact.
Represent $K$ by a complex $K^\bullet$ and consider the map
$$
K^\bullet
\longrightarrow
\bigoplus\nolimits_{n \geq 0} \tau_{\geq n} K^\bullet
$$
where we have used the canonical truncations, see
Homology, Section \ref{homology-section-truncations}.
This makes sense as in each degree the direct sum on the right is finite.
By assumption this map factors through a finite direct sum.
We conclude that $K \to \tau_{\geq n} K$ is zero for at least one $n$,
i.e., $K$ is in $D^{-}(R)$.

\medskip\noindent
Since $K \in D^{-}(R)$ and since every $R$-module is a quotient of a free
module, we may represent $K$ by a bounded above complex $K^\bullet$
of free $R$-modules, see
Derived Categories, Lemma \ref{derived-lemma-subcategory-left-resolution}.
Note that we have
$$
K^\bullet = \bigcup\nolimits_{n \leq 0} \sigma_{\geq n}K^\bullet
$$
where we have used the stupid truncations, see
Homology, Section \ref{homology-section-truncations}.
Hence by Lemma \ref{lemma-commutes-with-countable-sums} we see that
$1 : K^\bullet \to K^\bullet$ factors through
$\sigma_{\geq n}K^\bullet \to K^\bullet$ in $D(R)$.
Thus we see that $1 : K^\bullet \to K^\bullet$ factors as
$$
K^\bullet \xrightarrow{\varphi} L^\bullet \xrightarrow{\psi} K^\bullet
$$
in $D(R)$ for some complex $L^\bullet$ which is bounded and whose terms
are free $R$-modules. Say $L^i = 0$ for $i \not \in [a, b]$.
Fix $a, b$ from now on. Let $c$ be the largest integer $\leq b + 1$
such that we can find a factorization of $1_{K^\bullet}$ as above
with $L^i$ is finite free for $i < c$. We will show by induction that
$c = b + 1$. Namely, write $L^c = \bigoplus_{\lambda \in \Lambda} R$.
Since $L^{c - 1}$ is finite free we can find a finite subset
$\Lambda' \subset \Lambda$ such that $L^{c - 1} \to L^c$ factors
through $\bigoplus_{\lambda \in \Lambda'} R \subset L^c$. Consider the
map of complexes
$$
\pi :
L^\bullet
\longrightarrow
(\bigoplus\nolimits_{\lambda \in \Lambda \setminus \Lambda'} R)[-c]
$$
given by the projection onto the factors corresponding to
$\Lambda \setminus \Lambda'$ in degree $c$.
By our assumption on $K$ we see that, after possibly replacing $\Lambda'$ by
a larger finite subset, we may assume that $\pi \circ \varphi = 0$
in $D(R)$. Let $(L')^\bullet \subset L^\bullet$ be the kernel of $\pi$.
Since $\pi$ is surjective we get a short exact sequence of complexes,
which gives a distinguished triangle in $D(R)$ (see
Derived Categories, Lemma \ref{derived-lemma-derived-canonical-delta-functor}).
Since $\Hom_{D(R)}(K, -)$ is homological (see
Derived Categories, Lemma \ref{derived-lemma-representable-homological})
and $\pi \circ \varphi = 0$, we can find a morphism
$\varphi' : K^\bullet \to (L')^\bullet$ in $D(R)$ whose
composition with $(L')^\bullet \to L^\bullet$ gives $\varphi$.
Setting $\psi'$ equal to the composition of $\psi$ with
$(L')^\bullet \to L^\bullet$ we obtain a new factorization.
Since $(L')^\bullet$ agrees with $L^\bullet$ except in degree $c$
and since $(L')^c = \bigoplus_{\lambda \in \Lambda'} R$ the
induction step is proved.

\medskip\noindent
The conclusion of the discussion of the preceding paragraph is that
$1_K : K \to K$ factors as
$$
K \xrightarrow{\varphi} L \xrightarrow{\psi} K
$$
in $D(R)$ where $L$ can be represented by a finite
complex of free $R$-modules. In particular we see that $L$ is
perfect. Note that $e = \varphi \circ \psi \in \text{End}_{D(R)}(L)$
is an idempotent. By Derived Categories,
Lemma \ref{derived-lemma-projectors-have-images-triangulated}
we see that $L = \Ker(e) \oplus \Ker(1 - e)$.
The map $\varphi : K \to L$ induces an isomorphism with
$\Ker(1 - e)$ in $D(R)$. Hence we finally conclude that
$K$ is perfect by Lemma \ref{lemma-summands-perfect}.
\end{proof}